% page 90 of the facsimile (image p-089 of the scan)
% block 160.0 x 229.0 mm ; left margin 8.0 mm
\pgc{-12.700mm}{0.000mm}{
\l{\cel{7}82}
\l{}
\l{}
\l{\sou{celiba}. Qua vivas sen mariajir su. - DEFIS.}
\l{}
\l{\sou{celulo}.\cel{2}I. Mikra chambro (ex.: di monako, di karcerano).-}
\l{\cel{3}II.Alveolo quan konstruktas la abeli. - III.Mikra kavajo}
\l{\cel{3}en ula organi di la animali e di la vejetanti. - IV.}
\l{\cel{3}Elemento anatomiala embriona che la planti o la animali,}
\l{\cel{3}- DEFIRS.}
\l{}
\l{\sou{celluloido}. Kompozajo kemiala, qua imitas skalio, koralio,}
\l{\cel{3}e c. - DEFIRS.}
\l{}
\l{\sou{celuloso}. Korpo neutra ye qua konsistas la tisui di la veje-}
\l{\cel{3}tanti e specale la tisuo celulara ipsa. - DEFIRS.}
\l{}
\l{\sou{cembro}. (\sou{bot.}) Nomo di speco di pino, di qua la grani esas}
\l{\cel{3}manjebla. -L.\cel{1}\sou{cembro}. - F}
\l{}
\l{\sou{cemento}. Mixuro de kalko e de briki pistita, por interjuntar}
\l{\cel{3}la petro-bloki en masonuro. - DEFIRS.}
\l{}
\l{\sou{cementaca}r (\sou{trans.})\cel{2}Pozar stratope, sur metalo, substanco qua,}
\l{\cel{3}efike da temperaturo tre alta, modifikas lu, adportante a}
\l{\cel{3}lu un ek sua elementi. - DEFIRS.}
\l{}
\l{\sou{ceno}. I. Parto di teatro ube pleas la aktori. - II. Loko ube}
\l{\cel{3}eventas la pleajo. - III. Divido di akto teatrala, deter-}
\l{\cel{3}minita da la ekiro di la protagonisti. - DEFIRS.}
\l{}
\l{\sou{cenobit}o. En la komenco di la historio di la Eklezio Kristana,}
\l{\cel{3}la personulo qua vivis komune kun altra regulieri (opoze ad}
\l{\cel{3}\textquotedbl{}anakoreto\textquotedbl{}). - DEFIS.}
\l{}
\l{\sou{cenotafio}. Imituro di tombo, erektita memorige di mortinto,}
\l{\cel{3}e qua ne kontenas lua reliquii. - DEFIS.}
\l{}
\l{\sou{censo}. La revenuo evaluata da la autoritati statala, sive}
\l{\cel{3}por taxo, imposto, sive por ula yuro (elektiveso). - DEFIRS.}
\l{}
\l{\sou{censoro}. En Roma antiqua, un ek la du judiciisti di qui la}
\l{\cel{3}tasko konsistis ye okupar su pri la censo, reprimandar e}
\l{\cel{3}punisar la personi qui violacis la mori. - DEFIRS.}
\l{}
\l{\sou{censurar}\cel{2}(\sou{trans}.) I. Reprimandar ulu publike. - II. (\sou{aludante}}
\l{\cel{3}\sou{la Eklezio Kristana o korpo di ta Eklezio})\cel{2}Kondamnar}
\l{\cel{3}publike doktrino. - DEFIRS.}
\l{}
\l{\sou{cent}.\cel{3}lOO.\cel{2}- DEFIRS.}
\l{}
\l{\sou{centauro}. En la mitologio antiqua, ento mi-vira, mi-kavala.-}
\l{\cel{3}DEFIRS.}
}
